../cictikz/src/cictikz/data/tex/cictikz_preamble.tex

% Second step of the R-2R build-up: one series R in front of the section
% dac_2r_0 reduced to R, so looking in from V_1 the ladder is R + R = 2R.
% The current into the summing node is set by the tap voltage, I_0 = V_0/2R,
% and V_0 is half of V_1 because the R and the R it drives are equal.

\begin{circuitikz}[american, thick, transform shape, circuit ee IEC]

  ../cictikz/src/cictikz/data/tex/cictikz_lib.tex
  ../cictikz/src/cictikz/data/tex/rdac_lib.tex

  \def\yr{3.0}

  \coordinate (v1) at (0,\yr);
  \draw (v1) \hresistor{$R$} coordinate (v0);
  \node[anchor=south east] at ($(v1)+(-0.1,0.1)$) {$V_1$};
  \node[anchor=south] at ($(v0)+(0,0.15)$) {$V_0$};

  % ---- 2R across to the summing node ----
  \draw (v0) \hresistor{$2R$} -- ++(0.7,0) coordinate (vg);
  \rdacVgnd{(vg)}

  % ---- 2R from the tap down to ground ----
  \draw ($(v0)+(0,-1.6)$) \vresistor{$2R$};
  \draw ($(v0)+(0,-1.6)$) -- ++(0,-0.3) \vground;

  % ---- The current that leaves the tap, as the original marks it ----
  \draw[-{Latex[length=3mm]}] ($(v0)+(1.3,-0.6)$) -- ++(0,-1.0)
    node[anchor=west, midway, xshift=2mm] {$I_0$};

\end{circuitikz}
\end{document}
